%% TR Part C-ready manuscript skeleton — DREAM.
%% Compile with: pdflatex main.tex && bibtex main && pdflatex main.tex && pdflatex main.tex
%% This file is a SKELETON; numeric results, citations and figure files are inserted by
%% scripts/make_paper_figures.py + paper/narrative_report.md once M6/M7 complete.
\documentclass[3p,times,sort&compress,authoryear]{elsarticle}
\usepackage{amsmath,amssymb,amsfonts}
\usepackage{graphicx}
\usepackage{booktabs}
\usepackage{siunitx}
\usepackage[colorlinks=true,allcolors=blue]{hyperref}
\usepackage{xcolor}
\usepackage[caption=false]{subfig}
\sisetup{group-separator={\,}, output-decimal-marker={.}}

\journal{Transportation Research Part C}

\begin{document}

\begin{frontmatter}

\title{Operationalising ICAO Annex 14 Obstacle Limitation Surfaces into
Dynamic Safety Envelopes for eVTOL Airport Shuttle Integration:
Evidence from Real Airport Operational Data}

\author[institution]{Author Placeholder}
\address[institution]{Affiliation Placeholder}

\begin{abstract}
This study proposes \textbf{DREAM} (Dynamic Runway-configuration-aware Envelope for eVTOL
Airport Mobility), a real-data-driven dynamic safety-envelope framework for integrating
electric vertical-take-off-and-landing (eVTOL) airport-shuttle operations into conventional
airport environments. Instead of treating eVTOL airport access as a generic urban-air-mobility
routing problem, the study operationalises ICAO Annex~14 obstacle limitation surfaces (OLS)
into machine-readable three-dimensional constraints. Using Los Angeles International Airport
(LAX) as the primary case study and San Francisco International Airport (SFO) as an external
validation case, the framework integrates Federal Aviation Administration (FAA) airport
geometry, obstacle data, terrain elevation, OpenSky ADS-B trajectories, aviation weather
data, and airport ground-access demand data. A signed-distance-field representation is first
developed to encode static aerodrome protection surfaces. Runway-configuration-aware dynamic
envelopes are then generated by incorporating real runway-use patterns, wind conditions,
visibility, and conventional arrival/departure flows. A data-driven risk field is learned
from historical ADS-B trajectories through counterfactual injection to quantify spatiotemporal
interference risks between candidate eVTOL corridors and fixed-wing airport operations.
Finally, envelope-constrained corridor planning and simulation-based safety, capacity, and
accessibility assessment evaluate whether eVTOL shuttles can improve airport access without
compromising runway protection and airport capacity.
\end{abstract}

\begin{keyword}
eVTOL \sep airport access \sep ICAO Annex 14 \sep
obstacle limitation surfaces \sep dynamic safety envelope \sep
ADS-B \sep counterfactual injection \sep urban air mobility
\end{keyword}

\end{frontmatter}

%% --- Body ----------------------------------------------------------------------------
\section{Introduction}
\label{sec:intro}
% Motivation: eVTOL as airport access; gap: existing UAM literature treats it as generic LAM
% routing problem; airport-specific constraints (Annex 14, runway throughput) are usually
% absent. Contribution: DREAM — operationalising Annex 14 into a dynamic, runway-config-aware
% envelope, learning a real-data risk field, evaluating safety-capacity-accessibility jointly.
\textcolor{red}{[Filled by M8 from \texttt{paper/narrative\_report.md}.]}

\section{Related work}
\label{sec:related}
\subsection{eVTOL corridor and vertiport planning}
\subsection{Airport obstacle protection and Annex 14}
\subsection{Airport capacity assessment and ADS-B analytics}
\textcolor{red}{[Citations to be added; see manuscript bibliography.]}

\section{Methodology: DREAM}
\label{sec:method}
\subsection{Annex 14 as 3-D signed-distance constraints}\label{sec:sdf}
\subsection{Runway-configuration-aware dynamic envelope}\label{sec:envelope}
\subsection{ADS-B-driven risk-field learning}\label{sec:riskfield}
\subsection{Envelope-constrained A* corridor planner}\label{sec:planner}
\subsection{Safety--capacity--accessibility joint evaluation}\label{sec:eval}

\section{Data}
\label{sec:data}
% Eight real data sources; counterfactual-injection statement; primary window = Fri Aug 2024.
\subsection{Airport geometry, obstacles, terrain, ground network}
\subsection{Fixed-wing trajectories (OpenSky ADS-B)}
\subsection{Weather (NOAA AWC, ERA5)}
\subsection{Ground-access demand and passenger O\&D}

\section{LAX case study}
\label{sec:lax}
% B0–B4 across V1–V4 across 5 Fridays × 3 peak hours; figures from results/eval/KLAX.
\textcolor{red}{[Numbers + figures injected by \texttt{scripts/make\_paper\_figures.py}.]}

\section{SFO external validation}
\label{sec:sfo}

\section{Discussion}
\label{sec:discussion}
\subsection{Implications for regulation}
\subsection{Implications for airport planning}
\subsection{Threats to validity}
% (1) No real eVTOL data — counterfactual injection only.
% (2) Annex 14 numeric parameters are public-secondary-source — not normative ICAO text.
% (3) OpenSky ADS-B coverage below 1000 ft AGL is partial.

\section{Conclusion}
\label{sec:conclusion}

\section*{Reproducibility}
Code, data manifests, and configurations: \texttt{git@.../airportaccess}. Every figure and table
in this paper is regenerated by \texttt{bash scripts/run\_all.sh} on a fresh clone.

\section*{Acknowledgements}

\bibliographystyle{elsarticle-harv}
\bibliography{refs}

\end{document}
